\section{Reaching the edge}
\label{sec:reachingedge}

Section~\ref{sec:supplyestab} ended on a tension. The two factors we measured (how many R mutations each population produces, and what fraction of those mutations succeed at establishing in the edge) both differ between biofilm and well-mixed in roughly dose-independent ways: the active biofilm produces about $1.8\times$ as many mutations as the well-mixed reference at every dose, and each biofilm mutation establishes about $0.8\times$ as often at every dose. Multiplied together, these two ratios predict a constant biofilm-to-well-mixed rescue ratio. The actual rescue ratio is not constant. It climbs from below $1\times$ at $\dEdge = 1.05$ to about $2\times$ at $\dEdge = 2.8$. Some dose-dependent factor that those two measurements did not capture must therefore be carrying that flip. The candidate is the one piece of the rescue calculation that the two panels of Figure~\ref{fig:decomposition} do not cover: whether a core-born R mutation ever reaches the drug-exposed edge before going extinct. The supply panel counts every R mutation a population produces, but a core-born mutation contributes to rescue only if it arrives in the edge, where antibiotic selection makes resistance an advantage. The establishment panel measures what happens to mutations once they are in the edge, but not whether they got there. Whether a given core-born R lineage actually reaches the edge is therefore a separate per-lineage outcome from either of the two pieces Section~\ref{sec:supplyestab} measured, and we measure it directly across the same dose sweep.

Figure~\ref{fig:reaching} Panel A reports this fraction directly. Two curves are shown: the default active-core biofilm ($\bCore = 0.2$, open squares) and the more active core used in Section~\ref{sec:doubleedge}'s bar plot ($\bCore = 0.8$, filled triangles). Both curves rise sharply with dose. At $\dEdge = 1.05$, only about $5$ to $10$ percent of core-born R lineages ever reach the edge; at $\dEdge = 2.8$, that fraction climbs to between $40$ and $80$ percent depending on core activity. At every dose the less active core actually has a higher fraction of core lineages reaching the edge than the more active core, an inversion we return to in the mechanism discussion below. The main point is the dose-dependent shape: in both cases the fraction rises sharply between low and high dose, matching the rescue flip in Figure~\ref{fig:dosecurves}.

\begin{figure}[h!]
\centering
\includegraphics[width=0.95\textwidth]{delivery.png}
\caption{Per-lineage outcome of core-born R mutations. Panel A: fraction of core-born R lineages that ever reach the drug-exposed edge before going extinct, vs dose. Panel B: mean number of cells in a lineage when it first occupies an edge position, vs dose. Both panels show the default active core ($\bCore = 0.2$, open squares) and the more active core ($\bCore = 0.8$, filled triangles) used for Figure~\ref{fig:doubleedge}. Each point averages $100{,}000$ runs for $\bCore = 0.2$ and $5{,}000$ runs for $\bCore = 0.8$ per parameter combination. \emph{Parameters}: $\Nzero = 1000$, $l = 100$, $\mu = 10^{-4}$.}
\label{fig:reaching}
\end{figure}

The shape of Panel A follows from a race between two processes inside the core. A core-born R lineage starts as a single cell with balanced birth and death rates ($\bCore = \dCore$), so it does not grow on average; instead it fluctuates in size and, given enough time and turnover of cells in the core, is likely to go extinct by drift. The only way the lineage can matter for rescue is if the retreating boundary reaches its position before drift kills it. The boundary moves inward at speed $\sE \cdot l$ per unit time, where $\sE = \dEdge - \bEdge$ is the net edge death rate. At high dose, $\sE$ is large; the boundary moves inward quickly and overtakes core lineages before drift has had time to eliminate them. At low dose, $\sE$ is small; the boundary creeps inward very slowly, and most core lineages drift to extinction before they are ever exposed to the edge. The dose-dependence visible in Panel A is just this race playing out across the dose sweep, where the speed of population decline, matched with the relative size of the edge $l$, sets the time available for a core lineage to survive before the boundary catches it.

The same race explains a feature of Panel A that runs against naive expectation. At every dose, the less active core ($\bCore = 0.2$) has a higher fraction of core lineages reaching the edge than the more active core ($\bCore = 0.8$), even though the more active core produces roughly four times more core mutations per unit time. One might expect that more mutations would mean more chances for at least one to survive the race to the edge. But boundary speed and core drift are decoupled. Boundary retreat is set entirely by edge dynamics, so the time a lineage waits for the boundary to reach it is the same at $\bCore = 0.2$ as at $\bCore = 0.8$. The strength of drift in the core, however, depends directly on $\bCore$: a more active core means more births and deaths per unit time per cell, which means faster fluctuations in lineage size and faster extinction by chance. With $\bCore = 0.8$, lineages are killed in the core more quickly than at $\bCore = 0.2$, so a smaller fraction survive long enough for the boundary to reach them. The more active core produces more mutations in absolute terms, but it also kills them off faster, and the fraction surviving long enough to reach the edge drops with $\bCore$ even as the absolute mutation count rises with $\bCore$.

The dormant case is the limit of this mechanism at $\bCore = 0$. With no births in the core, no R lineages are ever born there in the first place; the race has nothing to enter. The boundary retreats at the same speed $\sE \cdot l$ as in the active case, but with no core mutations to catch, no core-born lineages ever reach the edge at any dose. The dormant biofilm therefore has no rescue advantage at any dose, not because the boundary fails to retreat, but because the core never produces a mutation for the boundary to expose to the edge in the first place.

Panel B reports a secondary consequence of the same race. The mean size of a core-born R lineage when it first arrives in the edge is slightly larger at low dose than at high dose: about $1.4$ cells per arrival at $\dEdge = 1.05$, dropping toward $1$ cell at higher dose. A lineage that survives a long wait in the core (which happens at low dose, when the boundary moves slowly) has had more time to fluctuate upward in size by chance before being caught. A lineage caught quickly (at high dose, when the boundary moves fast) is more likely to still be a single cell on first arrival. The dose dependence in Panel B runs opposite to what the rescue flip needs. A lineage arriving in the edge as a cluster of two or three cells has more independent chances to escape early extinction than a single-cell arrival: each cell can carry the lineage forward even if its siblings die. Larger clusters are therefore advantageous for establishment, and Panel B shows that larger clusters happen at low dose. If Panel B were the dominant influence on the rescue flip, the biofilm would rescue more often at low dose than at high dose. But the rescue flip goes the other way: the biofilm rescues less often at low dose and more often at high dose. Panel B therefore cannot be carrying the dose-dependence of the rescue flip.

What dominates instead is Panel A, and the reason is in the magnitudes. Panel A spans an order of magnitude across the dose range (from a few percent reaching the edge to about half reaching). Panel B varies by less than fifty percent (from about $1.4$ down to about $1.0$ cells per arrival). At low dose, the modest cluster advantage of a $1.4$-cell arrival cannot offset the fact that almost no core-born R lineages are arriving in the first place. The dose-dependence of the rescue flip therefore lives in the fraction reaching the edge at all, not in the size of those arrivals. Putting the three measurements together accounts for the rescue flip from Section~\ref{sec:doubleedge}. Rescue requires either an edge-born R mutation to establish (the path open to all populations) or a core-born R mutation to reach the edge and then establish (the path open only to biofilms with an active core). The first path is similar in biofilm and well-mixed: edge mutation counts overlap, and the per-mutation establishment fraction is only slightly reduced for the biofilm. The second path is what carries the dose-dependence.

At low dose, fewer than ten percent of core-born R lineages ever reach the edge. The active core produces more mutations overall, but those extra mutations stay trapped in the core where they cannot contribute to rescue. The biofilm therefore rescues less often than well-mixed at low dose, because the slight reduction in per-mutation establishment is no longer offset by anything from the core. At high dose, almost half of core-born R lineages reach the edge, so the active core's extra mutations now contribute a meaningful number of R lineages to the edge population, and the biofilm rescues more often than well-mixed. The crossover near $\dEdge \approx 1.4$ is where the dose-dependent reach-the-edge fraction balances the dose-independent reduction in per-mutation establishment.

The dormant case follows from the same accounting. With no core births, the second path is closed; only the first path exists, with the same edge mutation count as well-mixed and the same slightly reduced per-mutation establishment as the active biofilm. So the dormant biofilm rescues less often than well-mixed at every dose, exactly what Figure~\ref{fig:dosecurves} shows. The flip we set up in Section~\ref{sec:doubleedge} is the active biofilm's second path overtaking the first at high dose, and the dormant biofilm's lack of a second path at any dose. What remains to test is whether this account holds across the model parameters we have so far held fixed: the biofilm geometry $l$, the core activity $\bCore$, and the mutation rate $\mu$. The next section sweeps each in turn.

